\chapter{Introduction}

%How do we know what is grammatical?

%\textit{Pedantic}, meaning `plodding, out of date' from \textit{ped-} = `foot'. At least, that's what I thought until my mid twenties


%\section{The ungrammatical asterisk}

\begin{quote}
    Consider for a moment what grammar is. It is the most elementary part of logic. It is the beginning of the analysis of the thinking process. The principles and rules of grammar are the means by which the forms of language are made to correspond with the universal forms of thought\dots~The structure of every sentence is a lesson in logic. \hfill -- John Stuart Mill (1867:15)
\end{quote}

If you have any association at all with the asterisk $\langle$*$\rangle$, it's likely connected footnotes or perhaps the marking of birth dates. It's a symbol of additional information, a small detour from the main text. But within the world of linguistics, this little star has become the marker of ungrammaticality.

When I tell folks this, they imagine the linguist as a kind of teacher marking up student papers with asterisks. But that's not it at all. Often the ungrammatical sentences linguists are ``starring'' are those they've made up themselves. Rarely is it to call attention to the kinds of mistake that people make. In fact, it's often to illustrate the kinds of mistakes people \textit{do not} make. If this sounds odd, it is, and the starred sentences are regularly so wonderfully bizarre that I sometimes marvel at the ingenuity needed to come up with them. To give you a sense of what I mean, (\ref{ex:weird-ungrammaticals}) presents six typical oddballs from \textit{The Cambridge grammar of the English language}.

\ea \label{ex:weird-ungrammaticals}
    \ea[*]{\textit{He offered her the steak fiendishly hungry.}}
    \ex[*]{\textit{An apparently bird hit the car.}}
    \ex[*]{\textit{The news his guilt was devastating.}}
    \ex[*]{\textit{It had been for you.}}
    \ex[*]{\textit{She lives here any longer.}}
    \ex[*]{\textit{She wrote more plays than he wrote five novels.}}
    \z
\z

In this book, I will try to explain that feeling you get when you read these crazy examples, where it comes from, and what's special about it. This means you're going to look at lot of ungrammatical sentences and I'm going to use a lot of asterisks. So, it's worth asking: why an asterisk?

\bigskip

August Schleicher became a linguist -- or philologist -- before dictionaries could trace the history of our words with the precision we have today, before scholars confidently linked the English word \textit{father} with German \textit{Vater} and Sanskrit \textit{pitar}, for instance. Born in 1821, Schleicher was never content with the less systematic approaches common in traditional philology. His background in philosophy had fuelled a desire to question linguistic assumptions and seek scientific rigour. In an era when language study was often steeped in mysticism, Schleicher brought a naturalist's eye to the field. He saw words as evolving like biological organisms, subject to discoverable laws of change.

Traditional philologists of the time, influenced by theories like Grimm's Law, often focused on identifying regular patterns of sound shifts between languages. This approach, while groundbreaking, revealed the what of language change but left much of the why unanswered. Schleicher sought a more comprehensive system, one that could explain not just the regularities, but also why some changes caught on -- even to the point of taking over -- while many died out, in short, he wanted to understand linguistic evolution.

His approach was not merely academic. He was deeply influenced by the progressive ideals that spurred the student revolutions of 1848. These uprisings reflected a desire throughout much of Europe for greater freedoms and national self-determination, ideals that his father subtly endorsed and ideals which likely shaped young Schleicher's rejection of traditional authority. His father, a physician, had encouraged him to abandon a potential career in theology and instead immerse himself in the world of philosophy. It was here, under the tutelage of luminaries like Hegel, that Schleicher first encountered the captivating idea that seemingly disparate phenomena -- words, species, ideas -- might all be governed by the same underlying forces of development and change.

But philosophy, with its abstractions and hypotheticals, ultimately proved unsatisfying for Schleicher. He yearned for something tangible, something he could sink his teeth into. And he found it in the study of ancient languages.

\bigskip

For centuries, the history of the world's languages had had been clear only to those who placed deep faith my myths like the Tower of Babel. Among the more naturally minded, similarities between words were often dismissed as coincidences. In the late 18th century, Sir William Jones, a British judge stationed in India, changed all that while studying Sanskrit, the ancient sacred language of the subcontinent. Jones began to see a system in the striking parallels between Sanskrit, Greek, and Latin, in words for family members, body parts, and fundamental concepts that seemed to echo across millennia and continents. He theorized that these languages must descend from a common ancestor, a prehistoric root from which the diverse branches of a language family had grown.

Jones's idea was just the beginning, a seed needing fertile ground. August Schleicher -- alongside figures like Hensleigh Wedgwood (Darwin's brother-in-law), Frederick William Farrar, and Max Müller -- would cultivate this seed and transform language study from speculation to systematic analysis. Driven by an obsessive desire for clarity and an unshakable belief in language as a natural phenomenon, Schleicher aimed to impose order on the wilderness of words.

He approached languages with the precision of an anatomist, meticulously comparing words across a vast array of sources: musty old Germanic texts, Slavic folk songs, and the rigid structure of ancient Sanskrit. It wasn't simply about finding similar words but about dissecting them to uncover their underlying patterns.

Take the seemingly unremarkable English word \textit{tooth}. With a scholar's zeal, Schleicher traced its history, recognizing its cousin \textit{Zahn} in German or \textit{tand} in Dutch. Moving further back, he found \textit{tunþus} in Gothic, an extinct Germanic language. But he didn't stop there. Latin \textit{dens}, Greek \textit{odontos}, and Sanskrit \textit{danta}, these too he connected through precise analysis of sound shifts and grammatical patterns. Eventually, going beyond what anyone else had tried, he hypothesized their shared origin: a reconstructed word that nobody had spoken, read, or perhaps even imagined in thousands of years.

But there was a communicative challenge too. Schleicher needed a way to signal the missing links in the chain, the words that must have existed yet were lost to history. While there's written evidence for words like \textit{Zahn}, \textit{tunþus}, and \textit{danta}, he needed a way to mark the hypothetical, like *\textit{dent} from which they must have descended. He wasn't afraid to guess at its existence, but he wanted to be clear that it was only a guess. And so he hit on the idea of marking the unattested words with an asterisk.

This choice was a clever repurposing of an established scholarly symbol. In textual criticism, pioneered by figures like Aristophanes of Byzantium, the asterisk marked anomalies like repeated lines. Later, Origen used it to signal passages in the Old Testament present in some translations but absent from others \citep{Grafton2010}. This tradition of marking gaps, omissions, or uncertainties with a visually distinct symbol likely resonated with Schleicher. His reconstructed words were, in a sense, textual gaps within the larger story of language evolution, placeholders for forgotten sounds and forms.

Schleicher's approach to tracing word histories echoed the era's fascination with classification and genealogy. Naturalists were busy sorting beetles and birds into intricate family trees. Schleicher aimed to do the same with language, driven by the belief that words, like organisms, evolve according to predictable laws. Darwin's \textit{Origin of Species}, when Schleicher encountered it a few years after its release, resonated powerfully with his own linguistic work. After all, hadn't Darwin himself been inspired by languages, by their family trees and subtle shifts over time? In 1863, Schleicher boldly solidified this connection with a book probing the parallels between Darwinian evolution and the development of language. He embraced the radical notion: languages weren't artifacts that we had wrought but ``natural organisms which, beyond human will and in accord with definite laws, are born, grow, develop, age, and die; languages thus manifest that series of phenomena which we normally view as aspects of life. Glottics, or the science of language, is hence a natural science'' \citep[7]{schleicher1863} cited in \citep[82]{Goldsmith2019}.

%Schleicher's path was far from smooth. Traditional philologists, steeped in the study of literary texts and untroubled by questions of linguistic origins, saw his work as overly speculative, even frivolous. His groundbreaking ideas clashed with the conservative academic environment, earning him scorn and hindering his career advancement.

%Today, linguists see Schleicher as only partially on the right track. Language, while undeniably subject to historical change, isn't simply an organism. Darwinian evolution emphasizes slow, gradual modification through natural selection. But language is also shaped by cultural forces, wars, migrations, and technological innovation. Words can be borrowed or forcibly imposed, not just inherited. Schleicher's family tree implies a neat branching, yet languages mix, influence, and ``contaminate'' each other in ways difficult to capture on a simple diagram. Also, he had strong views on linguistic progress and degeneration rejected by most linguists today \citep{Goldsmith2019}.

In 1868, Schleicher died from tuberculosis at the age of 47, leaving his ambitious projects unfinished. His theories were refined and sometimes rejected by later linguists, but his asterisk endured. It became a silent testament to a profound shift -- a moment when linguistics embraced its identity as a natural science, a relentless search for patterns in the seemingly chaotic evolution of human language.

\bigskip

While August Schleicher is the father of the asterisk in linguistics, he never used it to mean `ungrammatical'. That innovation is due to Henry Sweet (1845--1912). Sweet was a multifaceted scholar, known for groundbreaking contributions to phonetics, Old English studies, and language teaching methodology. Like Schleicher, he was a man who challenged established norms, advocating for a rigorous scientific approach to language at a time when traditional philology reigned.

Sweet's rebellious spirit and unconventional approach to scholarship can be attributed, in part, to his unique background and educational journey. Unlike many British scholars of his era, who followed a well-trodden path through elite schools and traditional classical studies, Sweet left for a period of study at Heidelberg University in Germany, and that's where his passion for language was ignited.

%Sweet came to hold a deep conviction that spoken language held the key to understanding linguistic development. This focus on the primacy of speech set him apart from many contemporary philologists who primarily analyzed written texts. His emphasis on phonetics, the precise description of speech sounds, led to publications like A Handbook of Phonetics (1877) and The Sounds of English (1908), foundational works in the field.

%Sweet's interests weren't limited to pronunciation. He delved into Old English, authoring seminal works such as An Anglo-Saxon Reader (1876) and The History of Language (1900). In these, he brought his meticulous analysis not just to individual sounds but to the grammatical structures and historical development of language itself.

Keenly observant of the mismatch between traditional grammar rules and the realities of everyday speech, Sweet developed a descriptive approach to grammar. His \textit{A  new English grammar, logical and historical} moved away from prescriptive notions of right and wrong usage, ideas the Schleicher, for instance, had fully subscribed to. In the preface to \textit{A new English grammar: Logical and historical} in 1892, Sweet wrote,

\begin{quote}
    As my exposition claims to be scientific, I confine myself
to the statement and explanation of facts, without attempting
to settle the relative correctness of divergent usages. If an
`ungrammatical' expression such as \textit{it is me} is in general
use among educated people, I accept it as such, simply
adding that it is avoided in the literary language. \phantom{~ }\hfill(p. xi)
\end{quote}

With this focus on the language in use, Sweet repurposed Schleicher's symbol. Where Schleicher marked words not directly observed, Sweet used the asterisk to signal sentences like those in (\ref{ex:weird-ungrammaticals}) that wouldn't be observed, at least not in the English of a competent speaker. The first time he used it appears to have been in ``thus from \textit{these tall men} we can infer \textit{these men are tall}, but we cannot make \textit{some Englishmen} into *\textit{Englishmen are some}, or \textit{half the island} into *\textit{the island was half}'' \citep[14]{Sweet}.

Today, the asterisk is as likely to mark an ungrammatical sentence as a reconstructed word. This shift, pioneered by Henry Sweet, reflects the evolution of linguistics itself. Curiously, many linguists who use the asterisk daily may be unaware of its complex history, proof of how deeply this symbol is woven into the fabric of the field.

\bigskip

In 1973, Fred Householder (1973) pointed out that the asterisk hides a crucial fact: what counts as `ungrammatical' varies. Different speakers, or even the same speaker in different situations, might have varying judgments on some sentences. The asterisk presents grammaticality as a simple yes/no, when in reality it's more complex

\section{My interest in grammar}
I'm often jealous, and slightly in awe, of folks with strong and detailed biographical memories. I can recall, for example, the fantastic feeling I had watching the appearance, out of the fog, of a stunning Indonesian harbour with small rocky islands jutting out of a perfectly calm dark sea, not long after the sun had risen. I was working on an 18m yacht out of Darwin, Australia, but I have no idea which island it is that teases my memory, or what year. I don't remember where we were coming from or planning to go. I could probably work out the year, but most other details may be irretrievable: the captain's name or that of the yacht (though I'm strangely sure of its size). I remember a married couple was also among the crew, and perhaps one of my crew mates may have been named Chris. Or perhaps Steve.

The captain and I argued, I know, about his change of plans and I was consequently put off somewhere on or near Komodo island. I recall somehow explaining that I wished to get to another island, being put into a wooden fishing boat and taken out into the straight where a large ferry eventually appeared, taking me on board. And where did I go next? Bali? Flores? I certainly spent time on each island. But the sequence is lost to me. It's all very unclear, like in a dream, glimpses, scenes, thoughts, but no coherent narrative and few specifics. Much of it seems implausible.

And so it is with remembering when I became interested in grammar. It feels like it should be a pivotal moment, yet the details are frustratingly elusive.

\bigskip

I have memories of hours, whole afternoons, spent on the sixth floor of the Kinokuniya book store in Shinjuku. This was the English-language floor, where you'd find the English novels along with translations of Japanese novels. There was a section for language-teaching materials, both for learning Japanese and English, both of which I made extensive use of.

I remember being surprised at the number of people standing and reading manga, magazines, and books in the book shops and convenience stores. Somebody told me it was so normal they had a noun for it: \textit{tachiyomi} `standing reading'. It wasn't clear to me what the social boundaries of tachiyomi were, but it seemed to license my going into Kinokuniya, picking up any book in the Cambridge Applied Linguistics series (the ones with the red covers), and reading until my legs and back hurt from standing so long. There was no way I was going to be able to buy any of those tomes, which typically sold for something like ¥8,000 (about CDN\$80). But there were many such afternoons of reading about language. When did they start?

\bigskip

I arrived in Tokyo in late fall, 1993, to look for an English-teaching job. I had no training as an English teacher and little experience except for a few private lessons I taught when I was in Chiang Mai, perhaps a year earlier. So, that could have been a spark for my interest. That's also the year that Steven Pinker's \textit{The language instinct} came out. My copy says, ``First HarperPerennial edition published 1995,'' but it's the 15th printing which seems to have been in 1997 so I have these various touch points that help me to date my growing fascination with grammar. My Temple University Japan M.Ed. transcript shows that took Ken Schaffer's New Grammars course in the fall of 1999, by which time I was well and truly gripped.

It occurs to me now that my interest was probably sparked at least as much by the need to learn Japanese as by the need to teach English. I had arrived in Japan knowing literally nothing about the language, and I soon decided I was going to learn it faster and better than the other English speakers around me, a task at which I was only partially successful. While I focused on memorizing vocabulary and sentence patterns, the sheer difference between Japanese grammar and English grammar constantly nagged at me. Phrases were built differently, verb conjugations seemed baffling, and particles added layers of meaning I was struggling to grasp. This frustration presented me with a puzzle. Still, in those days in Japan, I was more interested the pedagogy and psychology of language teaching and learning. It would take something else to turn my attention more to grammar itself.

\bigskip

I moved back to Canada with my wife and toddler in 2003, just before the first global SARS outbreak. I was looking for language teaching work in Toronto, but the city, along with Hong Kong, was SARS central, and students were fleeing. I didn't have the requisite certification to teach in a high school, and the universities, far from hiring English-language teachers, were laying them off. After four months of searching with no replies, I was finally interviewed for a full-time position at Humber College, a position I still hold today.

I remember, not long after arriving, perhaps a few months later, perhaps a few years, a rising sense of confusion and embarrassment over FANBOYS, a mnemonic for the co-ordinating conjunctions in English (\textit{for}, \textit{and}, \textit{nor}, \textit{but}, \textit{or}, \textit{yet}, \& \textit{so}), confusion because I didn't understand why this idea seemed so oddly central in the Humber English-teaching eco-system and embarrassment because I'd never encountered it before.

I recall turning to the \textit{Cambridge grammar of the English language} to try to sort it out. I'd seen a copy years earlier at the Temple University library in Japan, but not been able to make much of it. Humber also had a copy, so I thought I'd give it another try. This time, I was able to make more sense of it and realized that I needn't have been embarrassed about FANBOYS. The experience led me to start my blog \textit{English, Jack}, where the first of roughly 800 posts, on Friday, July 28, 2006 (I only know because I checked) was entitled ``The myth of FANBOYS''. It also led me into a deep dive of the Cambridge grammar. It was around this time that I also encountered Geoff Pullum's insightful articles on the philosophy of grammar, which further piqued my interest in the foundational questions of what grammar really is and how we acquire it. This was it. I was fully hooked.

Over the next decade, I submerged myself in the study of English grammar -- figuring out how it works and writing about what I learned. But something bigger started to intrigue me: the puzzle of how we know which sentences ``sound right'' and which don't. And even more fundamentally, what is grammar, anyway? Is a grammatical mistake like a social faux pas, a violation of an unwritten etiquette, or is it more fundamental, a glitch in the way our brains process language? This book is my exploration of these questions, a journey to understand what grammar is, how we know it, and why it it matters.